% ═══════════════════════════════════════════════════════════════════════════
% UNTERRICHTSLANGENTWURF: Letzte Bucharbeit vor Assessment
% Fach: Spanisch (Spätbeginner, jahrgangsübergreifend)
% Kurs: Klasse 10, 11, 12 | 10 Schüler
% Lehrwerk: A_tope.com (Cornelsen)
% ═══════════════════════════════════════════════════════════════════════════

\documentclass[11pt, a4paper]{article}

% ═══════════════════════════════════════════════════════════════════════════
% PAKETE
% ═══════════════════════════════════════════════════════════════════════════

\usepackage[utf8]{inputenc}
\usepackage[T1]{fontenc}
\usepackage[ngerman]{babel}
\usepackage{helvet}
\renewcommand{\familydefault}{\sfdefault}

\usepackage[margin=2.5cm, left=3.5cm, top=2.5cm, bottom=2.5cm]{geometry}
\usepackage{setspace}
\onehalfspacing

\usepackage{enumitem}
\usepackage{tabularx}
\usepackage{booktabs}
\usepackage{array}
\usepackage{longtable}
\usepackage{multirow}

\usepackage[table]{xcolor}
\usepackage{amssymb}
\usepackage{tcolorbox}
\usepackage{fancyhdr}
\usepackage{lastpage}
\usepackage{titlesec}
\usepackage{parskip}

% Hyperlinks (für Inhaltsverzeichnis)
\usepackage[hidelinks]{hyperref}

% ═══════════════════════════════════════════════════════════════════════════
% FARBEN
% ═══════════════════════════════════════════════════════════════════════════

\definecolor{spanisch}{HTML}{c41e3a}
\definecolor{grau}{HTML}{666666}
\definecolor{hellgrau}{HTML}{f5f5f5}
\definecolor{success}{HTML}{2a7a4b}
\definecolor{warningorange}{HTML}{b35c00}

% ═══════════════════════════════════════════════════════════════════════════
% ÜBERSCHRIFTEN
% ═══════════════════════════════════════════════════════════════════════════

\titleformat{\section}
    {\normalfont\large\bfseries\color{spanisch}}
    {\thesection}{1em}{}

\titleformat{\subsection}
    {\normalfont\normalsize\bfseries}
    {\thesubsection}{1em}{}

\titlespacing*{\section}{0pt}{18pt}{10pt}
\titlespacing*{\subsection}{0pt}{12pt}{6pt}

% ═══════════════════════════════════════════════════════════════════════════
% KOPF- UND FUSSZEILE
% ═══════════════════════════════════════════════════════════════════════════

\pagestyle{fancy}
\fancyhf{}
\renewcommand{\headrulewidth}{0.4pt}
\renewcommand{\footrulewidth}{0pt}

\lhead{\small Langentwurf Tier 1 -- Spanisch Spätbeginner}
\rhead{\small DS 4: Letzte Bucharbeit vor Assessment}
\cfoot{\small Seite \thepage\ von \pageref{LastPage}}

% ═══════════════════════════════════════════════════════════════════════════
% BEFEHLE
% ═══════════════════════════════════════════════════════════════════════════

% Spanisch-Highlight
\newcommand{\es}[1]{\textcolor{spanisch}{\textit{#1}}}

% Info-Box
\newtcolorbox{infobox}[1][]{
    colback=hellgrau,
    colframe=spanisch,
    boxrule=0.8pt,
    arc=0pt,
    left=8pt, right=8pt, top=6pt, bottom=6pt,
    #1
}

% Score-Box
\newtcolorbox{scorebox}{
    colback=white,
    colframe=spanisch,
    boxrule=1pt,
    arc=0pt,
    left=10pt, right=10pt, top=8pt, bottom=8pt,
    fonttitle=\bfseries,
    title=Didaktik-Score
}

% ═══════════════════════════════════════════════════════════════════════════
% DOKUMENT
% ═══════════════════════════════════════════════════════════════════════════

\begin{document}

% ═══════════════════════════════════════════════════════════════════════════
% DECKBLATT
% ═══════════════════════════════════════════════════════════════════════════

\thispagestyle{empty}

\begin{center}
    \vspace*{1.5cm}

    {\Large\bfseries Unterrichtslangentwurf}

    \vspace{0.3cm}

    {\large Tier 1 -- Vollständig}

    \vspace{0.5cm}

    {\large Fach: Spanisch (Spätbeginner)}

    \vspace{2cm}

    {\LARGE\bfseries\color{spanisch} Letzte Bucharbeit vor Assessment}

    \vspace{0.3cm}

    {\large Punto final / Evaluación (Gruppe A) -- La vida laboral (Gruppe B)}

    \vspace{2.5cm}

    \begin{tabular}{rl}
        \textbf{Doppelstunde:} & 4 von 10 (Beobachtungssequenz) \\[6pt]
        \textbf{Kurs:} & Jahrgangsübergreifend (Kl. 10, 11, 12) \\[6pt]
        \textbf{Kursgröße:} & 10 Schüler \\[6pt]
        \textbf{Datum:} & Dienstag, 13.01.2026 \\[6pt]
        \textbf{Lehrwerk:} & A\_tope.com (Cornelsen) \\[6pt]
        \textbf{Seiten:} & S. 33-34 (Gruppe A) / S. 106-107 (Gruppe B) \\[6pt]
    \end{tabular}

    \vspace{2cm}

    \begin{tabular}{rl}
        \textbf{Lehrkraft:} & \underline{\hspace{5cm}} \\[6pt]
        \textbf{Raum:} & \underline{\hspace{5cm}} \\[6pt]
        \textbf{Zeit:} & \underline{\hspace{5cm}} (90 Minuten) \\[6pt]
    \end{tabular}

    \vfill

    \hrule
    \vspace{0.3cm}
    {\small\color{grau} Beobachtungssequenz Februar 2026 -- Assessment am Mo 19.01.2026 (DS 5)}

\end{center}

\newpage

% ═══════════════════════════════════════════════════════════════════════════
% INHALTSVERZEICHNIS
% ═══════════════════════════════════════════════════════════════════════════

\tableofcontents
\newpage

% ═══════════════════════════════════════════════════════════════════════════
% 1. BEDINGUNGSANALYSE
% ═══════════════════════════════════════════════════════════════════════════

\section{Bedingungsanalyse}

\subsection{Statistik der Lerngruppe}

Die Lerngruppe besteht aus 10 Schülerinnen und Schülern der Klassen 10, 11 und 12. Es handelt sich um einen jahrgangsübergreifenden Kurs für Spätbeginner (Spanisch als neu einsetzende Fremdsprache ab Klasse 10).

\begin{center}
\begin{tabular}{|l|c|c|c|}
    \hline
    \textbf{Klassenstufe} & \textbf{Anzahl} & \textbf{Lernjahr} & \textbf{GER-Niveau} \\
    \hline
    Klasse 10 & 6 & 1 & A1-A2 \\
    Klasse 11 & 2 & 2 & A2-B1 \\
    Klasse 12 & 2 & 3 & B1 \\
    \hline
\end{tabular}
\end{center}

Die Heterogenität erfordert eine konsequente Differenzierung in zwei Lerngruppen (Gruppe A: Lernjahr 1, Gruppe B: Lernjahr 2-3).

\subsection{Gruppeneinteilung}

\textbf{Gruppe A (7 Schüler):} Klasse 10 + L.K. (Kl. 11)
\begin{itemize}[leftmargin=2em]
    \item Lehrwerkseiten: S. 30-36 (Unidad 1-2)
    \item Themen: Mi barrio, Orte, ser/estar/hay
\end{itemize}

\textbf{Gruppe B (3 Schüler):} Klasse 11-12 (stärkere)
\begin{itemize}[leftmargin=2em]
    \item Lehrwerkseiten: S. 102-107 (Themenfeld Berufe)
    \item Themen: La vida laboral, Berufe, Konjunktionen
\end{itemize}

\subsection{Schülerprofile (Gruppe A + L.K.)}

\begin{center}
\small
\begin{tabularx}{\textwidth}{|l|c|X|X|X|}
    \hline
    \textbf{Kürzel} & \textbf{Niveau} & \textbf{Stärken} & \textbf{Schwächen} & \textbf{Strategie} \\
    \hline
    E.H. (Kl.10) & TOP & Franz. Hintergrund, 15 Punkte, immer aktiv & -- & Peer-Tutor, anspruchsvollere Aufgaben \\
    \hline
    L.N. (Kl.10) & stark & Proaktiv, versteht schnell, spricht gern & Etwas faul, abgelenkt & In Bewegung halten, Verantwortung geben \\
    \hline
    C.P. (Kl.10) & solide & Gute Aussprache, 12 Punkte, ruhig & Braucht Zeit zum Verstehen & Klare Anweisungen, Wartezeit \\
    \hline
    V.S. (Kl.10) & solide & Konstant, fleißig, ruhig & Weniger Selbstvertrauen & Kleine Erfolgserlebnisse, Lob \\
    \hline
    L.K. (Kl.11) & solide & Versteht gut, 10 Punkte & Ablenkbar, weniger fleißig & Struktur, klare Erwartungen \\
    \hline
    H.P. (Kl.10) & Nachzügler & Bemüht, 10 Punkte & 1-2 Monate später angefangen, Grundlagenlücken & Scaffolding, Partnering mit E.H. \\
    \hline
    A.A. (Kl.10) & unsicher & Fleißig, lernt viel, gute Testergebnisse & Unsicher, schaut zu E.H. & Ermutigung, getrennt von E.H. setzen \\
    \hline
    A.S. (Kl.10) & unsicher & Macht Fortschritte, 9 Punkte & Unsicher, Aussagen klingen wie Fragen & Positive Verstärkung, klare Strukturen \\
    \hline
\end{tabularx}
\end{center}

\subsection{Arbeitsvoraussetzungen}

\textbf{Materielle Voraussetzungen:}
\begin{itemize}[leftmargin=2em]
    \item Unterrichtsraum mit TV-Screen und iPad-Projektion (Apple TV)
    \item Kleine Kreidetafel für Gruppe B
    \item Lehrwerk A\_tope.com für alle Schüler vorhanden
    \item Arbeitsblätter werden analog ausgeteilt
\end{itemize}

\textbf{Methodische Voraussetzungen:}
\begin{itemize}[leftmargin=2em]
    \item Die SuS kennen die Arbeit in getrennten Gruppen (seit DS 1)
    \item Partnerarbeit und selbstständige Lehrbucharbeit sind etabliert
    \item Selbsttests (Evaluación) sind aus dem Lehrwerk bekannt
\end{itemize}

\subsection{Fachbezogener Entwicklungsstand}

\textbf{Gruppe A (vor dieser Stunde):}
\begin{itemize}[leftmargin=2em]
    \item Inhaltlich: Mi barrio (Stadtteile, Orte), ir a + Ort, poder/querer/tener que
    \item Grammatisch: ser/estar/hay, vuestro, Verbkonjugation (ir)
    \item Aussprache: [b]/[ß]-Unterscheidung
\end{itemize}

\textbf{Gruppe B (vor dieser Stunde):}
\begin{itemize}[leftmargin=2em]
    \item Inhaltlich: Spanisches Schulsystem, Arbeiten im Ausland (Texte Carla + Jan)
    \item Grammatisch: Konjunktionen (ya que, porque, mientras), Wortbildung
    \item Fertigkeiten: Mediación, Textverständnis
\end{itemize}

% ═══════════════════════════════════════════════════════════════════════════
% 2. SACHANALYSE
% ═══════════════════════════════════════════════════════════════════════════

\section{Sachanalyse}

\subsection{Gruppe A: Punto final + Evaluación (S. 33-34)}

\textbf{Punto final (S. 33):}

Die Aufgabe \glqq{}Punto final\grqq{} ist eine produktive Schreibaufgabe, bei der die SuS eine E-Mail an einen Austauschschüler verfassen. Diese integriert:
\begin{itemize}[leftmargin=2em]
    \item Sich vorstellen (Name, Alter, Wohnort)
    \item Den eigenen Stadtteil beschreiben (hay, es, está)
    \item Aktivitäten vorschlagen (ir a + Ort)
\end{itemize}

\textbf{Evaluación 1 (S. 34):}

Der Selbsttest überprüft die Kompetenzen aus Unidad 1-2:
\begin{itemize}[leftmargin=2em]
    \item Wortschatz: Orte, Aktivitäten, Stadtteile
    \item Grammatik: ser/estar/hay, Possessivpronomen, ir-Konjugation
    \item Fertigkeiten: Kurze Texte verstehen und produzieren
\end{itemize}

\textbf{Didaktischer Wert:}
Die Kombination aus produktiver Aufgabe (Punto final) und rezeptiver Überprüfung (Evaluación) bereitet optimal auf das mündliche Assessment vor, bei dem dieselben Strukturen im Partnergespräch verwendet werden.

\subsection{Gruppe B: La vida laboral (S. 106-107)}

\textbf{Aufgabe 3b (S. 106):}

Die SuS fassen zentrale Aussagen über das Arbeitsleben in Spanien zusammen, basierend auf dem zuvor gelesenen Text. Dies erfordert:
\begin{itemize}[leftmargin=2em]
    \item Leseverstehen: Informationen extrahieren
    \item Paraphrasieren: Mit eigenen Worten wiedergeben
    \item Konjunktionen: Begründungen formulieren (porque, ya que)
\end{itemize}

\textbf{Aufgabe 4 (S. 106):}

Interkulturelles Lernen: Die SuS erklären Unterschiede zwischen dem deutschen und spanischen Arbeitsleben. Aspekte:
\begin{itemize}[leftmargin=2em]
    \item Arbeitszeiten und Siesta-Tradition
    \item Befristete Verträge und Jugendarbeitslosigkeit
    \item Work-Life-Balance
\end{itemize}

\textbf{Aufgabe 5 (S. 107):}

Wortschatzarbeit: Synonyme finden für Begriffe aus dem Themenfeld Berufe/Bildung. Dies fördert:
\begin{itemize}[leftmargin=2em]
    \item Wortschatzerweiterung
    \item Sprachliche Flexibilität
    \item Vorbereitung auf freies Sprechen im Vortrag
\end{itemize}

\subsection{Assessment-Formate (Vorschau DS 5)}

\textbf{Gruppe A: Mündliche Partnerprüfung}
\begin{itemize}[leftmargin=2em]
    \item Thema: \glqq{}Mi barrio\grqq{} -- Partnergespräch
    \item Dauer: 3-4 Minuten pro Paar
    \item Bewertung: ser/estar/hay, Wortschatz, Grammatik, Flüssigkeit
\end{itemize}

\textbf{Gruppe B: Kurzvortrag}
\begin{itemize}[leftmargin=2em]
    \item Thema: \glqq{}Mi profesión ideal\grqq{} -- Einzelpräsentation
    \item Dauer: 3-4 Minuten pro Person
    \item Bewertung: Argumentation, Fachvokabular, Struktur, Flüssigkeit
\end{itemize}

% ═══════════════════════════════════════════════════════════════════════════
% 3. DIDAKTISCHE ANALYSE
% ═══════════════════════════════════════════════════════════════════════════

\section{Didaktische Analyse}

\subsection{Stellung in der Sequenz}

Dies ist die \textbf{4. Doppelstunde} der 10-DS-Beobachtungssequenz:

\begin{center}
\begin{tabular}{|c|l|l|}
    \hline
    \textbf{DS} & \textbf{Inhalt} & \textbf{Status} \\
    \hline
    1-2 & Bucharbeit parallel (A: S. 30-32, B: S. 102-104) & abgeschlossen \\
    3 & Bucharbeit + Assessment-Ankündigung & abgeschlossen \\
    \rowcolor{hellgrau}
    \textbf{4} & \textbf{Letzte Bucharbeit + Kriterien verteilen} & \textbf{heute} \\
    5 & \textbf{ASSESSMENT} (sonstige Note) & nächste Woche \\
    6-8 & Brückenthema + Konvergenz & Vorbereitung Beobachtung \\
    9 & Puffer & Reserve \\
    10 & \textbf{BEOBACHTUNG} (03.02.2026) & Ziel \\
    \hline
\end{tabular}
\end{center}

\subsection{Funktion der Stunde}

Diese Stunde hat eine \textbf{dreifache Funktion}:

\begin{enumerate}[leftmargin=2em]
    \item \textbf{Abschluss der Bucharbeit:} Letzte Seiten vor dem Assessment bearbeiten
    \item \textbf{Wiederholung:} Durch Evaluación (A) bzw. Zusammenfassung (B) Inhalte festigen
    \item \textbf{Assessment-Vorbereitung:} Bewertungskriterien transparent machen, Erwartungen klären
\end{enumerate}

\subsection{Kompetenzziele (Rahmenplan MV)}

\textbf{Kommunikative Kompetenz:}
\begin{itemize}[leftmargin=2em]
    \item Schreiben (A): Persönliche Texte verfassen (E-Mail)
    \item Sprechen (B): Zusammenhängend über Themen referieren
\end{itemize}

\textbf{Sprachliche Mittel:}
\begin{itemize}[leftmargin=2em]
    \item Grammatik (A): ser/estar/hay korrekt anwenden
    \item Grammatik (B): Konjunktionen zur Argumentation nutzen
\end{itemize}

\textbf{Interkulturelle Kompetenz:}
\begin{itemize}[leftmargin=2em]
    \item (B): Unterschiede im Arbeitsleben DE/ES reflektieren
\end{itemize}

\textbf{Methodische Kompetenz:}
\begin{itemize}[leftmargin=2em]
    \item Selbsteinschätzung durch Evaluación
    \item Bewertungskriterien verstehen und anwenden
\end{itemize}

\subsection{Legitimation}

\textbf{Bezug zum Rahmenplan Spanisch MV:}

Der Rahmenplan für die gymnasiale Oberstufe fordert die Entwicklung mündlicher und schriftlicher Produktionskompetenz. Die Vorbereitung auf mündliche Prüfungsformate (Partnergespräch, Vortrag) ist explizit vorgesehen.

\textbf{Bezug zur Bedingungsanalyse:}

Die heterogene Lerngruppe erfordert unterschiedliche Assessment-Formate. Das Partnergespräch (A) ermöglicht gegenseitige Unterstützung; der Einzelvortrag (B) entspricht dem fortgeschrittenen Niveau.

% ═══════════════════════════════════════════════════════════════════════════
% 4. STUNDENZIELE
% ═══════════════════════════════════════════════════════════════════════════

\section{Stundenziele}

\subsection{Grobziele}

\begin{enumerate}[leftmargin=2em]
    \item Die Schülerinnen und Schüler \textbf{festigen} ihre erworbenen Kompetenzen durch gezielte Wiederholung (Evaluación / Zusammenfassung).
    \item Die Schülerinnen und Schüler \textbf{kennen} die Bewertungskriterien für das Assessment und können sich gezielt vorbereiten.
\end{enumerate}

\subsection{Feinziele Gruppe A}

\begin{center}
\small
\begin{tabularx}{\textwidth}{|c|X|c|c|}
    \hline
    \textbf{Nr.} & \textbf{Die SuS...} & \textbf{AFB} & \textbf{Diff.} \\
    \hline
    A1 & ...verfassen eine strukturierte E-Mail an einen Austauschschüler unter Verwendung von ser/estar/hay. & II & alle \\
    \hline
    A2 & ...lösen die Aufgaben der Evaluación 1 mit mind. 70\% Korrektheit. & I-II & alle \\
    \hline
    A3 & ...benennen die vier Bewertungskriterien für das Partnergespräch. & I & alle \\
    \hline
    A4 & ...formulieren Beispielfragen für das Assessment (¿Hay...? ¿Cómo es...? ¿Adónde vas...?). & II & Standard \\
    \hline
\end{tabularx}
\end{center}

\subsection{Feinziele Gruppe B}

\begin{center}
\small
\begin{tabularx}{\textwidth}{|c|X|c|c|}
    \hline
    \textbf{Nr.} & \textbf{Die SuS...} & \textbf{AFB} & \textbf{Diff.} \\
    \hline
    B1 & ...fassen das Arbeitsleben in Spanien in eigenen Worten zusammen. & II & alle \\
    \hline
    B2 & ...erklären mind. 3 Unterschiede zwischen DE und ES im Bereich Arbeit. & II-III & alle \\
    \hline
    B3 & ...finden Synonyme für 5 vorgegebene Begriffe aus dem Themenfeld Berufe. & I & alle \\
    \hline
    B4 & ...benennen die vier Bewertungskriterien für den Kurzvortrag. & I & alle \\
    \hline
\end{tabularx}
\end{center}

% ═══════════════════════════════════════════════════════════════════════════
% 5. METHODISCHE ANALYSE
% ═══════════════════════════════════════════════════════════════════════════

\section{Methodische Analyse}

\subsection{Phasierung (90 Minuten)}

\begin{center}
\begin{tabular}{|c|l|c|}
    \hline
    \textbf{Phase} & \textbf{Inhalt} & \textbf{Zeit} \\
    \hline
    \multicolumn{3}{|c|}{\textbf{1. Stunde (45 Min)}} \\
    \hline
    Einstieg & Gemeinsam: Assessment-Ankündigung wiederholen & 5 Min \\
    \hline
    Arbeitsphase 1 & A: Punto final (E-Mail) / B: Aufg. 3b + 4 & 30 Min \\
    \hline
    Zwischensicherung & Kurzes Vorstellen erster Ergebnisse & 10 Min \\
    \hline
    \multicolumn{3}{|c|}{\textbf{2. Stunde (45 Min)}} \\
    \hline
    Arbeitsphase 2 & A: Evaluación 1 / B: Aufg. 5 + Vortrags-Gliederung & 20 Min \\
    \hline
    Kriterien & \textbf{Bewertungskriterien verteilen und besprechen} & 15 Min \\
    \hline
    Vorbereitung & Partnerbildung (A), Themenwahl (B) & 7 Min \\
    \hline
    Abschluss & Organisatorisches, Hausaufgabe & 3 Min \\
    \hline
\end{tabular}
\end{center}

\subsection{Sozialformen}

\begin{center}
\begin{tabular}{|l|l|l|}
    \hline
    \textbf{Phase} & \textbf{Sozialform} & \textbf{Begründung} \\
    \hline
    Einstieg & Plenum & Gemeinsamer Fokus auf Assessment \\
    \hline
    Punto final (A) & EA/PA & Individuelle Produktion, bei Bedarf Partnerhilfe \\
    \hline
    Aufg. 3b-4 (B) & EA & Konzentrierte Textarbeit \\
    \hline
    Evaluación (A) & EA & Selbsttest, individuelle Standortbestimmung \\
    \hline
    Aufg. 5 (B) & PA & Gemeinsame Wortschatzarbeit \\
    \hline
    Kriterien & Plenum (getrennt) & Gemeinsames Verständnis der Erwartungen \\
    \hline
    Vorbereitung & PA (A) / EA (B) & Assessment-spezifisch \\
    \hline
\end{tabular}
\end{center}

\subsection{Medien und Material}

\begin{center}
\begin{tabular}{|l|l|}
    \hline
    \textbf{Material} & \textbf{Verwendung} \\
    \hline
    Lehrwerk A\_tope.com & S. 33-34 (A), S. 106-107 (B) \\
    \hline
    Bewertungsbogen Gruppe A & Partnerprüfung (4 Kriterien, 18 Punkte) \\
    \hline
    Bewertungsbogen Gruppe B & Kurzvortrag (4 Kriterien, 18 Punkte) \\
    \hline
    TV-Screen & Projektion der Kriterien \\
    \hline
    Kleine Tafel & Notizen für Gruppe B \\
    \hline
    Evaluación-Lösungen (A) & Zur Selbstkontrolle \\
    \hline
\end{tabular}
\end{center}

\subsection{Differenzierung}

\textbf{Innerhalb Gruppe A:}
\begin{itemize}[leftmargin=2em]
    \item H.P., A.A., A.S. (unsicher): Vorformulierte Satzanfänge für E-Mail, mehr Zeit
    \item E.H., L.N. (stark): Erweiterte E-Mail mit Fragen an den Austauschschüler
\end{itemize}

\textbf{Innerhalb Gruppe B:}
\begin{itemize}[leftmargin=2em]
    \item Alle auf ähnlichem Niveau; Differenzierung durch Tiefe der Argumentation
\end{itemize}

\textbf{Schülerspezifische Maßnahmen:}
\begin{itemize}[leftmargin=2em]
    \item E.H. als Peer-Tutor für H.P. bei Fragen
    \item A.A. nicht neben E.H. setzen (Selbstständigkeit fördern)
    \item V.S. und C.P. als Partner (ähnliches Niveau, beide ruhig)
\end{itemize}

\subsection{Erwartete Schwierigkeiten}

\begin{center}
\begin{tabular}{|p{5cm}|p{8cm}|}
    \hline
    \textbf{Schwierigkeit} & \textbf{Reaktion} \\
    \hline
    A: Verwechslung ser/estar/hay & Wiederholung der Regeln an Tafel, Merkblatt \\
    \hline
    A: E-Mail-Format unbekannt & Beispiel-E-Mail zeigen (Punto final, S. 33) \\
    \hline
    B: Text zu komplex & Schlüsselwörter gemeinsam klären \\
    \hline
    Zeitproblem & Puffer: Aufg. 5 (B) kann als HA aufgegeben werden \\
    \hline
    Unsicherheit bei Kriterien & Beispiele geben, \glqq{}Was bedeutet Flüssigkeit?\grqq{} \\
    \hline
\end{tabular}
\end{center}

% ═══════════════════════════════════════════════════════════════════════════
% 6. VERLAUFSPLANUNG
% ═══════════════════════════════════════════════════════════════════════════

\section{Verlaufsplanung}

\subsection{1. Stunde (45 Minuten)}

\small
\begin{longtable}{|p{1.2cm}|p{1.8cm}|p{3.5cm}|p{3.5cm}|p{1.2cm}|p{1.5cm}|}
    \hline
    \textbf{Zeit} & \textbf{Phase} & \textbf{Lehrerhandeln} & \textbf{Schülerhandeln} & \textbf{SF} & \textbf{Medien} \\
    \hline
    \endfirsthead
    \hline
    \textbf{Zeit} & \textbf{Phase} & \textbf{Lehrerhandeln} & \textbf{Schülerhandeln} & \textbf{SF} & \textbf{Medien} \\
    \hline
    \endhead

    0-5 & Einstieg & Begrüßung. \glqq{}Erinnert ihr euch, was nächsten Montag ansteht?\grqq{} Assessment wiederholen. Überblick: Heute letzte Bucharbeit, dann Kriterien. & Antworten: \glqq{}Assessment\grqq{}, \glqq{}Prüfung\grqq{}. Hören zu, fragen nach. & PL & -- \\
    \hline

    5-10 & Überleitung & Gruppen trennen. A: S. 33, Punto final erklären. B: S. 106, Aufg. 3b-4 erklären. \glqq{}Ihr arbeitet selbstständig, ich wechsle zwischen den Gruppen.\grqq{} & Buch aufschlagen, Aufgaben lesen. & PL & Buch \\
    \hline

    10-40 & Arbeits\-phase 1 & \textbf{Bei A:} Hilfestellung bei E-Mail. Satzanfänge für Schwächere. Kontrolle ser/estar/hay. \newline \textbf{Bei B:} Textverständnis sichern, Schlüsselwörter klären. Diskussion anregen. & \textbf{A:} Schreiben E-Mail an Austauschschüler. Nutzen Vokabelliste S. 32. \newline \textbf{B:} Lesen Text, notieren Zusammenfassung. Vergleichen DE/ES. & EA/PA (A) EA (B) & Buch, Heft \\
    \hline

    40-45 & Zwischen\-sicherung & \textbf{A:} 1-2 SuS lesen E-Mail vor. Gemeinsame Korrektur. \newline \textbf{B:} Kurz: \glqq{}Was sind die Hauptunterschiede?\grqq{} & A: Vorlesen, zuhören, verbessern. \newline B: Nennen 2-3 Unterschiede. & PL (getrennt) & -- \\
    \hline
\end{longtable}
\normalsize

\subsection{2. Stunde (45 Minuten)}

\small
\begin{longtable}{|p{1.2cm}|p{1.8cm}|p{3.5cm}|p{3.5cm}|p{1.2cm}|p{1.5cm}|}
    \hline
    \textbf{Zeit} & \textbf{Phase} & \textbf{Lehrerhandeln} & \textbf{Schülerhandeln} & \textbf{SF} & \textbf{Medien} \\
    \hline
    \endfirsthead
    \hline
    \textbf{Zeit} & \textbf{Phase} & \textbf{Lehrerhandeln} & \textbf{Schülerhandeln} & \textbf{SF} & \textbf{Medien} \\
    \hline
    \endhead

    0-5 & Überleitung & Kurzer Rückblick 1. Stunde. \glqq{}Jetzt: Evaluación für A, Synonyme + Gliederung für B.\grqq{} & Hören zu, bereiten Material vor. & PL & -- \\
    \hline

    5-25 & Arbeits\-phase 2 & \textbf{A:} Evaluación S. 34 bearbeiten lassen. Danach Lösungen verteilen zur Selbstkontrolle. \newline \textbf{B:} Aufg. 5 (Synonyme). Dann: Gliederung für Vortrag \glqq{}Mi profesión ideal\grqq{} erstellen. & \textbf{A:} Lösen Evaluación, kontrollieren mit Lösung, markieren Fehler. \newline \textbf{B:} Finden Synonyme in PA. Erstellen Stichpunkte für Vortrag. & EA (A) PA/EA (B) & Buch, Lösung, Heft \\
    \hline

    25-40 & \cellcolor{hellgrau}\textbf{Kriterien} & \cellcolor{hellgrau}\textbf{Bewertungsbögen verteilen!} \newline Für jede Gruppe eigener Bogen. Kriterien erklären: \newline A: ser/estar/hay, Wortschatz, Grammatik, Flüssigkeit \newline B: Argumentation, Fachvokabular, Struktur, Flüssigkeit \newline \glqq{}Wer hat Fragen?\grqq{} Beispiele geben. & Lesen Kriterien. Fragen stellen: \glqq{}Was bedeutet ‚Kohärenz'?\grqq{} \glqq{}Wie lang genau?\grqq{} Notieren wichtige Punkte. & PL (getrennt) & Bewertungs\-bögen, TV \\
    \hline

    40-47 & Vorbereitung & \textbf{A:} Paare für Partnerprüfung festlegen. Paare üben kurz 2-3 Fragen. \newline \textbf{B:} Thema wählen: Welchen Beruf? Erste Stichpunkte. & A: Partnerzuordnung notieren, erste Probefragen stellen. \newline B: Berufswahl notieren, beginnen Argumentation. & PA (A) EA (B) & Heft \\
    \hline

    47-50 & Abschluss & HA: A vorbereitet kommen (Fragen üben), B Vortrag ausarbeiten. Ausblick: \glqq{}Montag ist es soweit!\grqq{} Positiver Zuspruch. & Notieren HA, räumen auf, verabschieden sich. & PL & -- \\
    \hline
\end{longtable}
\normalsize

\textbf{Legende:} PL = Plenum, EA = Einzelarbeit, PA = Partnerarbeit

% ═══════════════════════════════════════════════════════════════════════════
% 7. BEWERTUNGSKRITERIEN (ANLAGE)
% ═══════════════════════════════════════════════════════════════════════════

\section{Bewertungskriterien (Anlage)}

\subsection{Gruppe A: Mündliche Partnerprüfung \glqq{}Mi barrio\grqq{}}

\begin{center}
\begin{tabular}{|p{4.5cm}|c|c|}
    \hline
    \textbf{Kriterium} & \textbf{Punkte (0-3)} & \textbf{Gewichtung} \\
    \hline
    1. Verwendung ser/estar/hay & 0-3 & $\times$ 2 = max. 6 \\
    \hline
    2. Wortschatz (Stadtviertel, Orte) & 0-3 & $\times$ 1,5 = max. 4,5 \\
    \hline
    3. Grammatische Korrektheit & 0-3 & $\times$ 1 = max. 3 \\
    \hline
    4. Flüssigkeit + Verständlichkeit & 0-3 & $\times$ 1,5 = max. 4,5 \\
    \hline
    \multicolumn{2}{|r|}{\textbf{Gesamt:}} & \textbf{max. 18} \\
    \hline
\end{tabular}
\end{center}

\textbf{Ablauf:}
\begin{itemize}[leftmargin=2em]
    \item Partner A beschreibt sein Stadtviertel
    \item Partner B stellt 3 Fragen (¿Hay...? ¿Cómo es...? ¿Adónde vas...?)
    \item Rollentausch
    \item Dauer: 3-4 Minuten pro Paar
\end{itemize}

\textbf{Paare (nach Niveau gemischt):}
\begin{itemize}[leftmargin=2em]
    \item E.H. + A.S. (stark + unsicher)
    \item L.N. + A.A. (stark + unsicher)
    \item C.P. + V.S. (beide ruhig, ähnlich)
    \item L.K. + H.P. (beide Nachholbedarf)
\end{itemize}

\subsection{Gruppe B: Kurzvortrag \glqq{}Mi profesión ideal\grqq{}}

\begin{center}
\begin{tabular}{|p{4.5cm}|c|c|}
    \hline
    \textbf{Kriterium} & \textbf{Punkte (0-3)} & \textbf{Gewichtung} \\
    \hline
    1. Argumentation + Konjunktionen & 0-3 & $\times$ 2 = max. 6 \\
    \hline
    2. Fachvokabular (Berufe, Bildung) & 0-3 & $\times$ 1,5 = max. 4,5 \\
    \hline
    3. Struktur + Kohärenz & 0-3 & $\times$ 1 = max. 3 \\
    \hline
    4. Flüssigkeit + Aussprache & 0-3 & $\times$ 1,5 = max. 4,5 \\
    \hline
    \multicolumn{2}{|r|}{\textbf{Gesamt:}} & \textbf{max. 18} \\
    \hline
\end{tabular}
\end{center}

\textbf{Inhalt des Vortrags:}
\begin{enumerate}[leftmargin=2em]
    \item Welchen Beruf möchtest du? Warum?
    \item Welche Fähigkeiten brauchst du?
    \item Wie ist der Weg dorthin (Ausbildung/Studium)?
\end{enumerate}

\textbf{Format:} Freies Sprechen mit Stichpunkten, 3-4 Minuten

\subsection{Notenskala (14-Punkte-System)}

\begin{center}
\begin{tabular}{|c|c|c|c|c|c|c|}
    \hline
    18-17 & 16-15 & 14-13 & 12-11 & 10-9 & 8-7 & 6-5 \\
    \hline
    14 NP & 13 NP & 12 NP & 10 NP & 9 NP & 8 NP & 6 NP \\
    \hline
\end{tabular}
\end{center}

% ═══════════════════════════════════════════════════════════════════════════
% 8. HAUSAUFGABE
% ═══════════════════════════════════════════════════════════════════════════

\section{Hausaufgabe}

\textbf{Gruppe A:}
\begin{itemize}[leftmargin=2em]
    \item Vorbereitung auf Partnergespräch: 5 Fragen formulieren (¿Hay...? ¿Cómo es...? ¿Adónde vas...? etc.)
    \item Eigenen Stadtteil beschreiben können (mind. 6 Sätze)
    \item Optional: Mit Partner telefonisch üben
\end{itemize}

\textbf{Gruppe B:}
\begin{itemize}[leftmargin=2em]
    \item Vortrag \glqq{}Mi profesión ideal\grqq{} vollständig ausarbeiten (Stichpunkte)
    \item Mind. 3 Konjunktionen einbauen (ya que, porque, además, por eso)
    \item Vor dem Spiegel üben (3-4 Minuten)
\end{itemize}

% ═══════════════════════════════════════════════════════════════════════════
% 9. ANLAGEN
% ═══════════════════════════════════════════════════════════════════════════

\section{Anlagen}

\begin{enumerate}[leftmargin=2em]
    \item \textbf{Bewertungsbogen Gruppe A:} Partnerprüfung \glqq{}Mi barrio\grqq{} (1 Seite)
    \item \textbf{Bewertungsbogen Gruppe B:} Kurzvortrag \glqq{}Mi profesión ideal\grqq{} (1 Seite)
    \item \textbf{Lösungen Evaluación 1} (für Selbstkontrolle Gruppe A)
    \item \textbf{Beispiel-E-Mail} (Punto final, falls SuS unsicher)
    \item \textbf{Satzanfänge für schwächere SuS} (Hilfsblatt)
\end{enumerate}

\subsection*{Anlage 5: Satzanfänge für E-Mail (Gruppe A)}

\begin{infobox}
\textbf{Satzanfänge für deine E-Mail:}

\begin{itemize}[leftmargin=2em]
    \item \es{Hola, me llamo... y tengo... años.}
    \item \es{Vivo en..., que es/está...}
    \item \es{En mi barrio hay...}
    \item \es{Mi barrio es... porque...}
    \item \es{Cuando vienes, podemos ir a...}
    \item \es{¿Qué te gusta hacer?}
\end{itemize}
\end{infobox}

% ═══════════════════════════════════════════════════════════════════════════
% 10. DIDAKTIK-SCORE
% ═══════════════════════════════════════════════════════════════════════════

\newpage
\section{Didaktik-Score}

\begin{scorebox}
\textbf{Bewertung der didaktischen Qualität nach 10 Dimensionen}\\
\textbf{Ziel-Score: $\geq$ 9.0 (Premium-Qualität)}

\vspace{1em}
\begin{tabular}{|c|l|c|c|p{4.5cm}|}
    \hline
    \textbf{\#} & \textbf{Dimension} & \textbf{Gew.} & \textbf{Score} & \textbf{Notiz} \\
    \hline
    1 & Differenzierung & 15\% & 9/10 & Zwei Lerngruppen, Binnendiff. durch Partnering \\
    \hline
    2 & Taxonomie (AFB) & 10\% & 9/10 & I-III vertreten: Reproduktion, Transfer, Produktion \\
    \hline
    3 & Lernziel-Alignment & 10\% & 10/10 & Direkte Vorbereitung auf Assessment-Formate \\
    \hline
    4 & Aktivierung & 10\% & 9/10 & Produktive Aufgaben: E-Mail, Vortrag-Gliederung \\
    \hline
    5 & Scaffolding & 10\% & 9/10 & Satzanfänge, Peer-Tutoring, gestufte Hilfen \\
    \hline
    6 & Feedback-Qualität & 10\% & 8/10 & Selbstkontrolle Evaluación, aber wenig Prozess-Feedback \\
    \hline
    7 & Sprachliche Klarheit & 10\% & 9/10 & Kriterien explizit, Beispiele gegeben \\
    \hline
    8 & Motivation \& Relevanz & 10\% & 9/10 & Assessment-Bezug motiviert; Berufswahl relevant \\
    \hline
    9 & Inklusion (UDL) & 10\% & 8/10 & Verschiedene Zugänge (mündlich/schriftlich), könnte mehr visuell sein \\
    \hline
    10 & Praktikabilität & 5\% & 10/10 & 90 Min realistisch, Material vorhanden \\
    \hline
    \multicolumn{3}{|r|}{\textbf{GESAMT:}} & \textbf{9.0/10} & \textcolor{success}{\textbf{FREIGABE}} \\
    \hline
\end{tabular}

\vspace{1em}
\textbf{Bewertungsschwellen:}
\begin{itemize}
    \item[\checkmark] \textcolor{success}{$\geq$ 9.0: \textbf{FREIGABE} (Premium-Qualität)}
    \item[$\Box$] \textcolor{warningorange}{8.0--8.9: Iteration empfohlen}
    \item[$\Box$] 6.0--7.9: Überarbeitung erforderlich
    \item[$\Box$] $<$ 6.0: Grundlegende Neukonzeption
\end{itemize}
\end{scorebox}

\vfill

\hrule
\vspace{0.3cm}
\begin{center}
    {\small\color{grau} Langentwurf erstellt für Beobachtungssequenz Februar 2026 -- DS 4 von 10}
\end{center}

\end{document}
